\chapter{Motivation}
% Introduccion

\section{Uso y poder de los LLMS}

% Estadísticas de uso y aplicaciones. Que se note son heavys

\section{Pero qué es un LLM?}

% Muchos dibujos, diagramas de transformers
% Hablar de CoT, y de la aleatoriedad

\section{Formalizaciones previas y resultados que ignoran los aspectos probabilísticos}

% Intuición de los resultados
% Que se note que es más tardío lo de cot. Justificar por qué el COT es el parámetro relevante.
%   Contar que otras personas estudian parámetros como la profundidad, o embedding, etc.
% Remarcar que nadie modela los aspectos probabilísticos (salvo uno)

\section{Objetivos y contribuciones}

% Eso, qué esperamos encontrar.
% Preguntas concretas que motiven el trabajo, ej
%    -> Los resultados anteriores se extienden al contexto probabilistico?
%    -> Hay una separacion estilo P y BPP?
%    -> La clase probabilistica de CoT se comporta bien (clausura, amplificacion)
